\Needspace{18\baselineskip}
\section{Occurrence fibres and transport of hole observations}
The next milestone completes context closure under explicit hole freshness,
including replacement at arbitrary box depth. Its proofs need to distinguish
individual occurrences, rather than merely sets of labels. This section provides
coordinates for those occurrences without changing the existing hypernet type.

\Needspace{24\baselineskip}
\subsection{Coordinates local to an outer edge}
\begin{definition}[Children and occurrence fibre]\label{term:fibre}
For a label $l$ at depth zero, its child occurrence type is empty. At depth $d+1$,
a primitive label again has no children; a box label with body $B$ has child type
$\operatorname{AllEdge}(B)$. Define the fibre of $l$ to be
\[
 F_d(l)=\{\bullet\}+\operatorname{Children}_d(l).
\]
The left element is the edge itself; a right element is an occurrence below it.
\end{definition}
\lean{Graph/Occurrences}{12}{16}{Linear.Children}
\lean{Graph/Occurrences}{19}{20}{Linear.Fiber}
\termdef{fibre-lookup}
The shallow hole observation reads the primitive hole label at depth zero or in
the primitive summand at positive depth, and returns no hole on a box label.
The child observation eliminates empty child types and otherwise uses the
body's existing hole lookup. The fibre observation combines these functions:
left elements use shallow lookup and right elements use child lookup.
\lean{Graph/Occurrences}{22}{25}{Linear.shallowHole}
\lean{Graph/Occurrences}{27}{31}{Linear.childHole}
\lean{Graph/Occurrences}{33}{36}{Linear.fiberHole}
\termdef{occurrence-coordinates}
An occurrence of a graph $G$ in fibre coordinates is a dependent pair
\[
 (e,x),\qquad e\in G_E,\quad x\in F_d(\lambda_G(e)).
\]
The first component is always an outer edge, even when $x$ goes deep into a box.
This representation depends on edge labels and occurrence addresses, not on
vertex incidence or a particular enumeration of the carriers.
\lean{Graph/Occurrences}{39}{40}{Linear.Occurrences}

\begin{lemma}[Occurrence coordinates]
The existing occurrence type $\operatorname{AllEdge}(G)$ is bijective to the
fibre-coordinate occurrence type.
\end{lemma}
\begin{proof}
At depth zero, send an edge $e$ to $(e,\mathsf{left}(\bullet))$; the inverse
forgets that singleton component. A right component would belong to the empty
type and is impossible. At positive depth, send a shallow edge to the same
left-tagged pair, and a deep address $(e,a)$ to $(e,\mathsf{right}(a))$.
The inverse removes the fibre tag to reconstruct the corresponding shallow or
deep address. Both inverse laws follow by cases on the tags and uniqueness of the unit value.
\end{proof}
\lean{Graph/Occurrences}{42}{60}{Linear.occurrenceBijection}
Define a graph's occurrence observation at $(e,x)$ by applying the fibre
observation to $\lambda_G(e)$ and $x$.
\lean{Graph/Occurrences}{62}{63}{Linear.occurrenceHole}
The bijection preserves hole observations. At depth zero this is definitional.
For a positive-depth shallow occurrence, split the label into primitive or box
and compute both lookups. For a deep occurrence, both lookups read the same
body address: exposing the label cases makes the equality reflexive. The original
hierarchy lookup now uses a directly dependent reader, avoiding unnecessary
proof-dependent casts while retaining the same definition by label cases.
\lean{Graph/Occurrences}{65}{77}{Linear.occurrenceBijection_hole}

\begin{lemma}[The context condition in fibre coordinates]
A graph has unique holes exactly when any two fibre-coordinate occurrences
observing the same hole label are equal.
\end{lemma}
\begin{proof}
For the forward direction, apply unique holes to the inverse images of the two
occurrences. Observation preservation supplies their lookup hypotheses. Map
the resulting equality forwards and cancel the inverse pairs. Conversely, map
the two original addresses forwards, use fibre-coordinate uniqueness and
observation preservation, and conclude by injectivity of the bijection.
\end{proof}
\lean{Graph/Occurrences}{80}{96}{Linear.uniqueHoles_iff_occurrences}

\Needspace{24\baselineskip}
\subsection{Transport along equality and edge maps}
If labels $l,l'$ are equal, transport of $x\in F_d(l)$ along that equality gives
an element of $F_d(l')$ with the same hole observation. Such transport is injective.
For each assertion, eliminate the label equality; transport becomes the identity,
so the observation equation is reflexive and injectivity is the given equality
of transported elements.
\lean{Graph/Occurrences}{98}{101}{Linear.fiberHole_transport}
\lean{Graph/Occurrences}{103}{106}{Linear.fiber_transport_injective}
A shallow edge map $f:G_E\to H_E$ with equations
$\lambda_G(e)=\lambda_H(f(e))$ maps all occurrences by
\[
 (e,x)\longmapsto (f(e),\operatorname{transport}(x)).
\]
The descendant address is preserved through the label equality; it is not
searched for again by comparing box contents.
\lean{Graph/Occurrences}{109}{112}{Linear.mapOccurrences}
This map preserves hole observation by the equality-transport observation lemma.
If $f$ is injective, the occurrence map is injective: equality of mapped pairs
first identifies their edge indices by injectivity of $f$. After substituting
that equality, equality of their transported second components implies equality
of the original components by transport injectivity. The dependent pairs are
therefore equal.
\lean{Graph/Occurrences}{114}{119}{Linear.mapOccurrences_hole}
\lean{Graph/Occurrences}{121}{132}{Linear.mapOccurrences_injective}
